\subsection{外尔不等式}

在本节中，我们考虑 K = C 上的希尔伯特空间。

设 E 是希尔伯特空间。对每个 \(n \in N\) ，我们回顾希尔伯特外幂 \(\hat{\Lambda}^{n}E\) 已在 EVT, V, p. 34 中定义。对每个希尔伯特空间 F 以及每个 \(u \in \mathcal{L}(E;F)\) ，我们还定义了线性映射 \(\hat{\Lambda}^{n}u \in \mathcal{L}(\hat{\Lambda}^{n}E;\hat{\Lambda}^{n}F)\) （同前）。这些构造是函子性的：对每个希尔伯特空间 G 以及每个线性映射 \(v \in \mathcal{L}(F;G)\) ，公式

\[
\widehat {\Lambda} ^ {n} v \circ \widehat {\Lambda} ^ {n} u = \widehat {\Lambda} ^ {n} (v \circ u)
\]

成立（同前，公式 (28)）。

设 H 是 E 的闭子空间，并设 \(i_{H}\) 是从 H 到 E 中的典范单射。对 E 的每个自同态 u，我们用 \(u_{H}\) 表示 H 的自同态 \(i_{H}^{*}ui_{H}\) 。

引理 2. --- 设 \(n \in N\) 。映射 \(\hat{\Lambda}^{n}i_{H}\) 是从 \(\hat{\Lambda}^{n}H\) 到 \(\hat{\Lambda}^{n}E\) 中的线性等距映射。

设 \((e_{j})_{j\in\mathrm{J}}\) 是 H 的标准正交基，\((e_{j})_{j\in\mathrm{J}^{\prime}}\) 是 E 的标准正交基，其中 \(J\subset J^{\prime}\) 。给 \(J^{\prime}\) 赋予一个全序。元素 \(e_{j_{1}}\wedge\cdots\wedge e_{j_{n}}\) （其中 \(j_{1}<\cdots<j_{n}\) 在 \(J^{\prime}\) 中（相应地，在 J 中）变化）构成 \(\widehat{\Lambda}^{n}E\) （相应地，\(\widehat{\Lambda}^{n}H\) ）的标准正交基（EVT, V, p. 34 的命题 5）。引理得证。

在下面，我们将把 \(\hat{\Lambda}^n\mathrm{H}\) 等同于 \(\hat{\Lambda}^n\mathrm{E}\) 的一个闭子空间，借助的是映射 \(\hat{\Lambda}^n i_{\mathrm{H}}\) 。

引理 3. --- 设 F 是希尔伯特空间，且 \(u \in \mathcal{L}(\mathrm{E};\mathrm{F})\)。设 \(n \in \mathbf{N}\) 。a) 我们有 ( \(\widehat{\Lambda}^n u\) ) \(^*\) = \(\widehat{\Lambda}^n (u^*)\) ；

\begin{enumerate}
\def\labelenumi{\alph{enumi})}
\setcounter{enumi}{1}
\item
  若 F = E 且 u 是埃尔米特的（相应地，正的、正规的、酉的），则 \(\hat{\Lambda}^{n}u\) 是埃尔米特的（相应地，正的、正规的、酉的）；
\item
  我们有 \(|\widehat{\Lambda}^n u| = \widehat{\Lambda}^n |u|\) ；
\item
  设 H 是 E 的闭子空间。限制 \((\widehat{\Lambda}^{n}u)|\widehat{\Lambda}^{n}\mathrm{H}\) （即 \(\widehat{\Lambda}^{n}u\) 在 \(\widehat{\Lambda}^{n}\mathrm{H}\) 上的限制）等于 \(\widehat{\Lambda}^{n}(u|\mathrm{H})\) （等式在 \(\mathcal{L}(\widehat{\Lambda}^{n}\mathrm{H};\widehat{\Lambda}^{n}\mathrm{F})\) 中成立）；
\item
  假设 F = E。设 H 是 E 的闭子空间。我们有 \((\hat{\Lambda}^{n}u)_{\hat{\Lambda}^{n}\mathrm{H}}=\hat{\Lambda}^{n}(u_{\mathrm{H}})\) （在 \(\mathcal{L}(\hat{\Lambda}^{n}\mathrm{H})\) 中）。
\end{enumerate}

断言 \(a\)）由 EVT, V, p. 39 的公式 (13) 得出。它立即蕴含：\(\widehat{\Lambda}^{n}u\) 在 \(u\) 是埃尔米特的（相应地，酉的）时是埃尔米特的（相应地，酉的）。若 \(u\) 是正的，则 \(u^{1/2}\) 是埃尔米特的，因而 \(\widehat{\Lambda}^{n}(u^{1/2})\) 也是；关系式 \(\widehat{\Lambda}^{n}u = (\widehat{\Lambda}^{n}u^{1/2})^{2}\) 蕴含 \(\widehat{\Lambda}^{n}u\) 是正的（I, p. 138，命题 8）。这就证明了 \(b\)）。

利用前述结果，可以算得

\[
\begin{array}{r l} & {\left(\widehat {\Lambda} ^ {n} | u |\right) ^ {2} = \widehat {\Lambda} ^ {n} (| u | ^ {2}) = \widehat {\Lambda} ^ {n} (u ^ {*} u)} \\ & {\qquad = \widehat {\Lambda} ^ {n} u ^ {*} \widehat {\Lambda} ^ {n} u = \left(\widehat {\Lambda} ^ {n} u\right) ^ {*} (\widehat {\Lambda} ^ {n} u) = | \widehat {\Lambda} ^ {n} u | ^ {2}.} \end{array}
\]

由于 \(\widehat{\Lambda}^n |u|\) 依 \(b\)）是正的，故由 I, p. 118 的命题 16 得 \(\widehat{\Lambda}^n |u| = |\widehat{\Lambda}^n u|\) ，由此得 \(c\)）。

设 \(i_{\mathrm{H}}\) 是从 H 到 E 中的典范单射。空间 \(\widehat{\Lambda}^{n}\mathrm{H}\) 等同于 \(\widehat{\Lambda}^{n}\mathrm{E}\) 的一个闭子空间，借助的是映射 \(\widehat{\Lambda}^{n}i_{\mathrm{H}}\) ，由此得

\[
(\widehat {\Lambda} ^ {n} u) | \widehat {\Lambda} ^ {n} \mathrm{H} = \widehat {\Lambda} ^ {n} u \circ \widehat {\Lambda} ^ {n} i _ {\mathrm{H}} = \widehat {\Lambda} ^ {n} (u \circ i _ {\mathrm{H}}) = \widehat {\Lambda} ^ {n} (u | \mathrm{H}),
\]

以及

\[
(\widehat {\Lambda} ^ {n} u) _ {\widehat {\Lambda} ^ {n} \mathrm{H}} = (\widehat {\Lambda} ^ {n} i _ {\mathrm{H}}) ^ {*} \circ \widehat {\Lambda} ^ {n} u \circ \widehat {\Lambda} ^ {n} i _ {\mathrm{H}} = \widehat {\Lambda} ^ {n} (i _ {\mathrm{H}} ^ {*} u i _ {\mathrm{H}}) = \widehat {\Lambda} ^ {n} (u _ {\mathrm{H}}),
\]

这就证明了 \(d\)）与 \(e\)）。

设 F 是希尔伯特空间，且 \(u \in \mathcal{L}^{c}(\mathrm{E};\mathrm{F})\) 。如同 IV, p. 151 的 no. 3 中那样，我们用 \(I_{E}\) 表示 E 的满足 \(F \neq E\) 的有限维子空间 F 的维数所成的集。我们回顾，\((\alpha_{n}(u))_{n \in I_{E}}\) 表示 u 的奇异值的扩张序列（IV, p. 156 的定义 3）。

命题 10. --- 我们有等式

\[
\prod_ {i = 0} ^ {n} \alpha_ {i} (u) = \left\| \widehat {\Lambda} ^ {n + 1} u \right\|
\]

对所有 \(n \in I_{E}\) 。

引理 3 的 c) 表明 \(\|\widehat{\Lambda}^{n+1}u\|=\|\widehat{\Lambda}^{n+1}|u\|\) ；此外，由于 \(\alpha_{i}(u)=\alpha_{i}(|u|)\) 对所有 \(i\in I_{E}\) 成立（IV, p. 157 的注 5, (2)），只需对 \(|u|\) 证明断言即可。于是可设 F=E 且 u 是正的。

于是设 \(\mathrm{B}=(e_{j})_{j\in\mathrm{J}}\) 是 E 的标准正交基，使得 u 在基 B 中可对角化（IV, p. 149 的定理 1），并设 \((\lambda_{j})_{j\in\mathrm{J}}\) 是 u 在基 B 中的特征值族。给 J 赋予一个全序，并用 \(J_{n}\) 表示 J 的元素的严格递增族 \((j_{0},\ldots,j_{n})\in\mathbf{J}^{n+1}\) 所成的集。向量

\[
e _ {\iota} = e _ {j _ {0}} \wedge \dots \wedge e _ {j _ {n}}
\]

（其中 \(\iota=(j_{0},\ldots,j_{n})\in\mathrm{J}_{n}\)）构成标准正交基 \(B_{n}\) ，它是 \(\widehat{\Lambda}^{n+1}\mathrm{E}\) 的标准正交基（EVT, V, p. 34，命题 5）。对每个 \(\iota=(j_{0},\ldots,j_{n})\in\mathrm{J}_{n}\) ，记

\[
\lambda_ {\iota} = \prod_ {j = 0} ^ {n} \lambda_ {i _ {j}}.
\]

由此可得 \((\hat{\Lambda}^{n + 1}u)e_{\iota} = \lambda_{\iota}e_{\iota}\) ，因而 \(\hat{\Lambda}^{n + 1}u\) 在基 \(\mathrm{B}_n\) 中是对角的。于是，我们有

\[
\| \widehat {\Lambda} ^ {n + 1} u \| = \sup _ {\iota \in \mathrm{I} _ {n}} \lambda_ {\iota}
\]

（IV, p. 147 的引理 1），它等于 \(\lambda_{0}(u)\cdots\lambda_{n}(u)\) ，即最大的 \(n+1\) 个特征值（\(u\) 的）之积。由于 \(\lambda_{i}(u)=\alpha_{i}(u)\) 对所有 \(i\in\mathrm{I}_{\mathrm{E}}\) 在 \(u\) 为正时成立（IV, p. 157 的注 5, (3)），即得所欲证的公式。

特别地，若 \(\alpha_{1}(u) < \alpha_{0}(u) = \| u\|\) ，则可看出，不等式 \(\| \widehat{\Lambda}^{n}(u)\| \leqslant \| u\|^{n}\) （EVT, V, p. 34，公式 (29)）在 \(n \geqslant 2\) 时一般不取等号。

推论. --- 设 G 是希尔伯特空间，且 \(v \in \mathcal{L}^{\mathrm{c}}(\mathrm{F};\mathrm{G})\) 。我们有

\[
\prod_ {i = 0} ^ {n} \alpha_ {i} (v \circ u) \leqslant \prod_ {i = 0} ^ {n} \alpha_ {i} (u) \alpha_ {i} (v)
\]

对所有 \(n \in \mathrm{I}_{\mathrm{E}}\) 。

只需注意

\[
\| \widehat {\Lambda} ^ {n + 1} (v u) \| = \| \widehat {\Lambda} ^ {n + 1} v \circ \widehat {\Lambda} ^ {n + 1} u \| \leqslant \| \widehat {\Lambda} ^ {n + 1} v \| \| \widehat {\Lambda} ^ {n + 1} u \|
\]

并应用命题 10。

引理 4. --- 设 A 是环。设 \(n \in \mathbf{N}\) ，并设 \((a_i)_{0 \leqslant i \leqslant n}\) 与 \((b_i)_{0 \leqslant i \leqslant n}\) 是 A 的元素族。对 \(0 \leqslant j \leqslant n\) ，置

\[
\mathrm{A} _ {j} = \sum_ {i = 0} ^ {j} a _ {i}.
\]

我们有

\[
\sum_ {i = 0} ^ {n} a _ {i} b _ {i} = \mathrm{A} _ {n} b _ {n} - \sum_ {i = 0} ^ {n - 1} \mathrm{A} _ {i} (b _ {i + 1} - b _ {i}).
\]

置 \(A_{-1}=0\) 。我们有 \(a_{i}=A_{i}-A_{i-1}\) （\(0\leqslant i\leqslant n\)），因而

\[
\sum_ {i = 0} ^ {n} a _ {i} b _ {i} = \sum_ {i = 0} ^ {n} (\mathrm{A} _ {i} - \mathrm{A} _ {i - 1}) b _ {i} = \sum_ {i = 0} ^ {n - 1} \mathrm{A} _ {i} (b _ {i} - b _ {i + 1}) + \mathrm{A} _ {n} b _ {n},
\]

即得所欲证。

设 I 是 \(\overline{R}\) 的一个不含 \(+\infty\) 的区间。回顾（FVR, I, p. 38，注），称连续函数 \(f: I \to [-\infty, +\infty[\) 是凸的，如果它在 I 的内部上的限制是取实值的凸函数；此时它在 inf I 处有（有限或无穷的）右极限。在下面的陈述中，按约定，0 与 \(\{-\infty, +\infty\}\) 中元素之积等于 0。

引理 5. --- 设 I ⊂ R 是一个不含 +∞ 的区间，f 是从 I 到 \([-\infty, +\infty[\) 中的一个递增凸函数。

设 n 是自然数，且

\[
a _ {0} \geqslant a _ {1} \geqslant \dots \geqslant a _ {n}, \quad b _ {0} \geqslant b _ {1} \geqslant \dots \geqslant b _ {n}
\]

是 I 的元素。

设 \(\varrho_{0} \geqslant \varrho_{1} \geqslant \cdots \geqslant \varrho_{n}\) 是正实数。假设

\[
\sum_ {i = 0} ^ {j} a _ {i} \leqslant \sum_ {i = 0} ^ {j} b _ {i}.\tag{5}
\]

对所有满足 \(0 \leqslant j \leqslant n\) 的整数 j 成立。则我们有

\[
\sum_ {i = 0} ^ {n} \varrho_ {i} f (a _ {i}) \leqslant \sum_ {i = 0} ^ {n} \varrho_ {i} f (b _ {i}).\tag{6}
\]

首先假设 \(a_i \in \mathring{\mathrm{I}}\) （从而 \(f(a_i) \in \mathbf{R}\) ）对所有 \(i\) 成立。设 \(i\) 满足 \(0 \leqslant i \leqslant n\) ，并设 \((\alpha_i, \beta_i)\) 是实数，使得方程为 \(y = \alpha_i x + \beta_i\) 的直线是 \(f\) 的图形在点 \((a_i, f(a_i))\) 处的支撑直线（FVR, I, p. 37）。于是有 \(f(a_i) = \alpha_i a_i + \beta_i\) 与 \(f(b_i) \geqslant \alpha_i b_i + \beta_i\) ，因为 \(f\) 的图形位于支撑直线的上方，由此得

\[
f (b _ {i}) - f (a _ {i}) \geqslant \alpha_ {i} (b _ {i} - a _ {i}).\tag{7}
\]

此外，若 \(i < n\) ，我们有

\[
\alpha_ {i} \geqslant f _ {g} ^ {\prime} (a _ {i}) \geqslant f _ {d} ^ {\prime} (a _ {i + 1}) \geqslant \alpha_ {i + 1}
\]

（同前及 FVR, I, p. 36，推论 1）；此外 \(\alpha_{i} \geqslant 0\) ，因为 \(f\) 递增（FVR, I, p. 22，推论），从而 \(\varrho_{i}\alpha_{i} \geqslant \varrho_{i+1}\alpha_{i+1} \geqslant 0\) （\(0 \leqslant i < n\)）。

设 \(j\) 是满足 \(0 \leqslant j \leqslant n\) 的整数。置

\[
\mathrm{A} _ {j} = \sum_ {i = 0} ^ {j} (b _ {i} - a _ {i}),
\]

于是由假设 (5) 知 \(A_{j} \geqslant 0\) 。依次应用不等式 (7) 与引理 4，我们导出

\[
\begin{array}{l} \sum_ {i = 0} ^ {n} \varrho_ {i} (f (b _ {i}) - f (a _ {i})) \geqslant \sum_ {i = 0} ^ {n} \varrho_ {i} \alpha_ {i} (b _ {i} - a _ {i}) \\ \qquad = \varrho_ {n} \alpha_ {n} A _ {n} + \sum_ {j = 0} ^ {n - 1} (\varrho_ {j} \alpha_ {j} - \varrho_ {j + 1} \alpha_ {j + 1}) A _ {j} \geqslant 0. \end{array}
\]

考虑一般情形，并对 \(n\) 作归纳推理。若某个 \(a_{i}\) 不属于 I 的内部，则必有 \(a_{0} = \sup I\) 或 \(a_{n} = \inf I\) 。

先设 \(a_{n}=\inf I\) 。此时有 \(a_{n}\leqslant b_{n}\) ，从而 \(\varrho_{n}f(a_{n})\leqslant\varrho_{n}f(b_{n})\) 。把归纳假设应用于族 \((a_{0},\ldots,a_{n-1})\) 、\((b_{0},\ldots,b_{n-1})\) 与 \((\varrho_{0},\ldots,\varrho_{n-1})\) ，便得

\[
\sum_ {i = 0} ^ {n - 1} \varrho_ {i} f (a _ {i}) \leqslant \sum_ {i = 0} ^ {n - 1} \varrho_ {i} f (b _ {i}),
\]

再加上 \(\varrho_{n}f(a_{n})\)即得所欲证的不等式。\(a_{0} = \sup I\) 的情形可类似处理。

命题 11（外尔不等式）. --- 设 G 是希尔伯特空间，且 \(v \in \mathcal{L}^{\mathrm{c}}(\mathrm{F};\mathrm{G})\)。设 \(g: \mathbf{R}_{+} \to [-\infty, +\infty[\) 是不减函数，使得函数 \(g \circ \exp\) 是凸的。我们有

\[
\sum_ {i = 0} ^ {n} g \left(\alpha_ {i} (v \circ u)\right) \leqslant \sum_ {i = 0} ^ {n} g \left(\alpha_ {i} (v) \alpha_ {i} (u)\right)
\]

对所有 \(n \in I_{E} \cap I_{F}\) 。

置 \(I = [-\infty, +\infty[\) 且 \(f = g \circ \exp\)。可以对下列取值应用引理 5

\[
a _ {i} = \log (\alpha_ {i} (v \circ u)) \in \mathrm{I}, \quad b _ {i} = \log (\alpha_ {i} (v) \alpha_ {i} (u)) \in \mathrm{I}
\]

并取 \(\varrho_{i} = 1\) （\(0\leqslant i\leqslant n\)），因为

\[
\sum_ {i = 0} ^ {j} a _ {i} = \log \Bigl (\prod_ {i = 0} ^ {j} \alpha_ {i} (v \circ u) \Bigr) \leqslant \log \Bigl (\prod_ {i = 0} ^ {j} \alpha_ {i} (v) \alpha_ {i} (u) \Bigr) = \sum_ {i = 0} ^ {j} b _ {i}
\]

对所有 \(0 \leqslant j \leqslant n\) 成立，依据是命题 10 的推论。于是不等式 (6) 就是所要的结论。

引理 6. --- 设 g 与 h 分别是定义在 R 的区间 I 与 J 上的凸函数。若 g 递增且定义在 h 的像上，则函数 g ∘ h 在 J 上是凸的。

事实上，对 \(t \in [0,1]\) 与 \((x,y) \in \mathrm{J} \times \mathrm{J}\) ，我们有

\[
\begin{array}{c} g (h (t x + (1 - t) y)) \leqslant g (t h (x) + (1 - t) h (y)) \\ \leqslant t g (h (x)) + (1 - t) g (h (y)). \end{array}
\]

推论. --- 设 G 是希尔伯特空间，且 \(v \in \mathcal{L}^{\mathrm{c}}(\mathrm{F};\mathrm{G})\)。设 \(n \in I_{E} \cap I_{F}\) 。

\begin{enumerate}
\def\labelenumi{\alph{enumi})}
\tightlist
\item
  设 \(r \in \mathbf{R}_{+}^{*}\) 。我们有
\end{enumerate}

\[
\sum_ {i = 0} ^ {n} \alpha_ {i} (v \circ u) ^ {r} \leqslant \sum_ {i = 0} ^ {n} \alpha_ {i} (v) ^ {r} \alpha_ {i} (u) ^ {r}.
\]

\begin{enumerate}
\def\labelenumi{\alph{enumi})}
\setcounter{enumi}{1}
\tightlist
\item
  设 \(p, q, r \in \mathbf{R}_{+}^{*}\) 满足 \(\frac{1}{p} + \frac{1}{q} = \frac{1}{r}\) 。则
\end{enumerate}

\[
\left(\sum_ {i = 0} ^ {n} \alpha_ {i} (v \circ u) ^ {r}\right) ^ {1 / r} \leqslant \left(\sum_ {i = 0} ^ {n} \alpha_ {i} (v) ^ {p}\right) ^ {1 / p} \left(\sum_ {i = 0} ^ {n} \alpha_ {i} (u) ^ {q}\right) ^ {1 / q}.
\]

\begin{enumerate}
\def\labelenumi{\alph{enumi})}
\setcounter{enumi}{2}
\tightlist
\item
  假设 F = E。对每个整数 \(m \geqslant 2\) ，我们有
\end{enumerate}

\[
\sum_ {i = 0} ^ {n} \alpha_ {i} (u ^ {m}) \leqslant \left(\sum_ {i = 0} ^ {n} \alpha_ {i} (u) ^ {2}\right) ^ {m / 2}.
\]

设 \(r \in R_{+}^{*}\) ，并设 g 是在 \(\mathbf{R}_{+}\) 上由 \(g(x) = x^{r}\) 定义的函数。函数 \(g \circ \exp\) 是凸的（引理 6），故把命题 11 应用于函数 g 即得断言 a)。

断言 \(b\)）由 \(a\)）及赫尔德不等式（INT, I，命题 4）得出。

设 \(m \geqslant 2\) 是整数。我们把 \(b)\) 应用于 \(r = 1\) 、\(p = q = 2\) 与 \(v = u^{m-1}\) ，得到

\[
\sum_ {i = 0} ^ {n} \alpha_ {i} (u ^ {m}) \leqslant \left(\sum_ {i = 0} ^ {n} \alpha_ {i} (u ^ {m - 1}) ^ {2}\right) ^ {1 / 2} \left(\sum_ {i = 0} ^ {n} \alpha_ {i} (u) ^ {2}\right) ^ {1 / 2}.
\]

我们对 \(m \geqslant 2\) 作归纳来证明 c)。前一不等式给出了 m = 2 时的断言。设 \(m \geqslant 3\) 且断言对 m - 1 成立；由于我们有

\[
\sum_ {i = 0} ^ {n} \alpha_ {i} (u ^ {m - 1}) ^ {2} \leqslant \left(\sum_ {i = 0} ^ {m} \alpha_ {i} (u ^ {m - 1})\right) ^ {2},
\]

上面的不等式蕴含

\[
\sum_ {i = 0} ^ {n} \alpha_ {i} (u ^ {m}) \leqslant \left(\sum_ {i = 0} ^ {n} \alpha_ {i} (u ^ {m - 1})\right) ^ {1 / 2} \left(\sum_ {i = 0} ^ {n} \alpha_ {i} (u) ^ {2}\right) ^ {1 / 2} \leqslant \left(\sum_ {i = 0} ^ {n} \alpha_ {i} (u) ^ {2}\right) ^ {m / 2}
\]

其中应用了归纳假设。

